% Abstract. Written last; every claim is checked against Sections 6 and 7 and is no stronger.
% Claims: r-v-h2-headline, r-h2-emp-h2-att/lo95/hi95, r-v-h2 (label probabilities),
% r-v-s1, r-v-s2, r-sec-cemn-s1-att, r-sec-dube-s2-att, r-v-h1, r-h1-wage-h1-att/t/c-hi95,
% r-tnu-h1-95, r-ph-h1-joint, r-v-h3, r-plc-nonff-sign.
\begin{abstract}
California's AB 1228 set a \$20 minimum wage for workers at large fast-food chains from April
2024, and the published QCEW-based estimates of its employment effect disagree. We report a
pre-specified, not blind, update for private limited-service restaurants (NAICS 722513) in
2024Q2--2026Q1, a window that adds the preliminary 2026Q1. A rule fixed before any post-period
outcome was computed picked the primary estimator from 24 candidates on data through 2019: a
trend-adjusted difference in differences with population-weighted donors from all other states.
Inference is built for one treated state. The pooled effect on log employment is $\HTwoAtt$
(95\% interval $\HTwoLoNinetyFive$ to $\HTwoHiNinetyFive$). The pre-specified label is
``inconclusive'', and it is qualified as not separable from a California-wide shock, because
employment outside 722513 also ran above its counterfactual. The design could not have shown an
effect smaller than 2\%: before any post-period data were read, the power run gave that label
probability $\ProbTauZeroUndSmaller$ even with no effect, so the result does not discriminate
between the published estimates. Re-estimations of both published approaches are also
inconclusive ($\SOneAtt$ and $\STwoAtt$). The weekly-earnings first stage, $\HOneAtt$ log
points, is not supported at the pre-specified critical value ($t = \HOneT$ against $\HOneCrit$),
which is not evidence that earnings did not rise; by the design's rule the employment estimate is
then uninformative about the wage floor. Two post hoc checks of that
verdict point in opposite directions. The year-two comparison is uninformative by rule.
\end{abstract}
